%%%%%%%%%%  scan idx 471  |  folio 460  %%%%%%%%%%

Para probar el lema, observemos que
\begin{equation*}
K_n=\ker\widetilde u_n
\end{equation*}
es también el anulador izquierdo del acoplamiento canónico
\begin{equation*}
u_n:\underline M_n\times\underline M'_n
 \longrightarrow(\mathbb Z/n\mathbb Z)_S,
\end{equation*}
cuya definición explícita se da en 9.4 y 9.6 en términos de las
extensiones
\begin{equation}\tag{11.6.14}
\resizebox{\textwidth}{!}{%
\begin{tikzcd}[ampersand replacement=\&, column sep=small, row sep=large]
0 \arrow[r]
 \& {\ensuremath{\check{\underline M}'_n(1)}} \arrow[r]
 \& F_n \arrow[r] \arrow[d, equal]
 \& R_n \arrow[r]
 \& 0
 \& {,}
 \& 0 \arrow[r]
 \& {\ensuremath{\check{\underline M}_n(1)}} \arrow[r]
 \& F'_n \arrow[r] \arrow[d, equal]
 \& R'_n \arrow[r]
 \& 0 \\
{}
 \& {}
 \& ({}_nA^o)^f
 \& {}
 \& {}
 \& {}
 \& {}
 \& {}
 \& ({}_nA'^o)^f
 \& {}
 \& {}
\end{tikzcd}}
\end{equation}
(donde $R_n$ y $R'_n$ son duales de Cartier entre sí), y en términos
de la extensión mixta $X_{n\eta}$ de $E_{n\eta}$ por $F_{n\eta}$,
fibra genérica de \sisi{(11.7.5)}{(11.6.5)}, donde ponemos
\begin{equation}\tag{11.6.15}
\begin{tikzcd}[column sep=large, row sep=large]
E_n=D(\sisi{F_n}{F'_n}),
 & 0 \arrow[r]
 & R_n \arrow[r] \arrow[d, no head, "\wr"']
 & E_n \arrow[r]
 & \underline M_n \arrow[r]
 & 0
 & {.} \\
{}
 & {}
 & D(R'_n)
 & {}
 & {}
 & {}
 & {}
\end{tikzcd}
\end{equation}
Se deduce inmediatamente de la proposición
\sisi{9.4.6}{9.4.7} que $K_n$ es el mayor de los subgrupos finitos
étale $K'_n$ de $\underline M_n$ tales que la imagen inversa
$\overline X_{n\eta}$ de $K'_n$ en la extensión mixta $X_{n\eta}$,
considerada como extensión mixta de $\overline E_{n\eta}$ por
$F_{n\eta}$, se prolonga a una extensión mixta de $\overline E_n$ por
$F_n$. Aquí $\overline E_n$ es la extensión evidente del subgrupo
\sisi{$K_n$}{$K'_n$} por $R_n$.

Esta condición se cumple evidentemente para el subgrupo ${}_n\Phi$ de
$\underline M_n$ en (11.6.11), pues se dispone de la correspondiente
extensión mixta $({}_nA)^f$ sobre $S$. Por tanto,
${}_n\Phi\subset K_n$. Para probar la inclusión opuesta, consideremos
la

%%%%%%%%%%  scan idx 472  |  folio 461  %%%%%%%%%%

extensión mixta $\overline X_n$ de $\overline E_n$ por $F_n$ que
prolonga $\overline X_{n\eta}$, definida como antes tomando
$K'_n=K_n$. Es, pues, una extensión

\begin{equation*}
0\longrightarrow F_n\longrightarrow\overline X_n
 \longrightarrow K_n\longrightarrow0,
\end{equation*}
y existe un homomorfismo canónico
\begin{equation}\tag{*}
F_n=({}_nA^o)^f\lhook\joinrel\longrightarrow A,
\end{equation}
que sobre la fibra genérica está inducido por el homomorfismo canónico
\begin{equation}\tag{**}
\overline X_{n\eta}\lhook\joinrel\longrightarrow X_{n\eta}
 ={}_nA_K\lhook\joinrel\longrightarrow A_K.
\end{equation}
Se deduce inmediatamente de (5.9.2) que el homomorfismo $(**)$ se
prolonga de manera única a un homomorfismo
\begin{equation}\tag{***}
\overline X_n\longrightarrow A
\end{equation}
que induce $(*)$. Como este último homomorfismo aplica la fibra genérica
en ${}_nA$, que es un subesquema cerrado de $A$, se factoriza a través
de $\overline X_n\longrightarrow{}_nA$; y, puesto que $\overline X_n$
es finito, se factoriza a través de $({}_nA)^f$. Examinando el
homomorfismo inducido sobre las fibras genéricas, concluimos que
\begin{equation*}
\overline X_{n\eta}\subset({}_nA)^f_\eta,
\quad\text{es decir,}\quad
K_{n\eta}\subset{}_n\Phi_\eta,
\quad\text{es decir,}\quad
K_n\subset{}_n\Phi.
\end{equation*}
Esto prueba el lema.

\medskip
\noindent
\underline{Observaciones 11.7.} La construcción realizada en
\sisi{11.7}{11.6} muestra que, a partir de las extensiones (11.6.14)
sobre $S$ y de la extensión mixta $X_{n\eta}={}_nA_K$ de $E_n$ por
$F_n$, se puede reconstruir, salvo isomorfismo único, la extensión
mixta (11.6.5) y, por tanto, también el subgrupo
\begin{equation*}
({}_nA)^f\supset({}_nA^o)^f
\end{equation*}
de ${}_nA$. En efecto, en virtud de 9.4, estos datos determinan el
acoplamiento
\begin{equation*}
u_n:\underline M_n\times\underline M'_n
 \longrightarrow(\mathbb Z/n\mathbb Z)_S,
\end{equation*}
y la construcción

%%%%%%%%%%  scan idx 473  |  folio 462  %%%%%%%%%%

precedente proporciona $({}_nA)^f$ como la extensión
mixta $\overline X_n$ de $\overline E_n$ por $F_n$ que prolonga la
extensión mixta $\overline X_{n\eta}$ inducida por $X_{n\eta}$. Aquí
$\overline E_n$ es la imagen inversa de $\ker\widetilde u_n$ por
$E_n\longrightarrow\underline M_n$. Por \sisi{(9.4.8)}{(9.4.7)}, esta extensión mixta
$\overline X_n$ queda así determinada \underline{salvo isomorfismo
único}. Se recupera entonces ${}_nA$ pegando $\overline X_n$ y
$X_{n\eta}$ a lo largo de su subesquema abierto común
$\overline X_{n\eta}$.

Por tanto, si se conocen los datos estructurales (9.6.1) y la extensión
mixta
\begin{equation*}
X_\eta=T_\ell(A_K)
\end{equation*}
de $E_\eta$ por $F_\eta$ ---equivalentemente, si se conocen las
extensiones (11.6.14) y $X_{n\eta}$ para todo
$n=\ell^{r+1}$, $r\geq0$---, se puede reconstruir el sistema inductivo
(o, si se prefiere, proyectivo) de extensiones mixtas (11.6.5).
